\chapter{Lokale Flachheit, kovariante Ableitung und Krümmung}

\lectureoverview{An einem ausgewählten Punkt lässt sich die Metrik stets bis zur ersten Ordnung euklidisch machen. Die Krümmung beginnt dort, wo diese Koordinatenfreiheit endet: beim Scheitern der Anpassung benachbarter lokaler Bezugssysteme in zweiter Ordnung. Wir entwickeln diese Aussage in drei Schritten. Zunächst betrachten wir Normalkoordinaten, dann die kovariante Ableitung und ihre Verbindung und schließlich die Nichtvertauschbarkeit zweier kovarianter Ableitungen, die den Riemann-Tensor definiert und der Gezeitenwirkung der Gravitation eine invariante Bedeutung gibt.}

\section{Die unvermeidliche Mühe mit den Indizes}

Die allgemeine Relativitätstheorie ist auf eine ganz bestimmte Weise schwierig. Ihre Grundideen sind nicht unverständlich, doch eine Rechnung ist schnell voller Indizes, Ableitungen und Symbole. Es gibt stärker komprimierte Sprachen: Vielbeine, Differentialformen, Spinoren, Twistoren und verwandte Hilfsmittel können lange Komponentengleichungen in kompakte Objekte verpacken. Diese Sprachen zu lernen ist jedoch selbst eine erhebliche Aufgabe, und kompakte Notation kann gerade die Rechenschritte verbergen, die man beim ersten Kennenlernen am besten verstehen sollte.

Die Maxwell-Gleichungen sind ein gutes Beispiel. Komponentenweise hingeschrieben bestehen sie aus mehreren einzelnen Gleichungen. Die Vektornotation fasst sie zu Divergenz- und Rotationgleichungen zusammen und macht so die Struktur leichter erkennbar. Auch in der allgemeinen Relativitätstheorie sind entsprechende Vereinfachungen möglich, doch wenn eine konkrete Rechnung tatsächlich ausgeführt werden muss, bleibt die Komponentenmethode der verlässliche Weg. Deshalb lassen wir die Indizes sichtbar. Sobald ihre Grammatik vertraut ist, wird der gedankliche Weg viel leichter, auch wenn die Rechnungen weiterhin seitenlang sein können.

Symbolische Computeralgebraprogramme können bei diesen Rechnungen helfen. Aus einer Metrik \(g_{mn}(x)\) kann ein solches Programm Christoffelsymbole, Riemann-Tensoren, Ricci-Tensoren und Einstein-Tensoren berechnen. Anschließend kann es diese Ausdrücke in die Feldgleichungen einsetzen. Das macht es nicht überflüssig, die Bedeutung der Objekte zu verstehen. Es trennt lediglich die mechanische Ausarbeitung von der Geometrie.

\begin{classroomqa}
\audiencequestion{Welche Forschungsgebiete waren zur Zeit dieser Vorlesung in der allgemeinen Relativitätstheorie besonders aktiv?}
\lecturerresponse{In der Vorlesung wurden mehrere Forschungsrichtungen hervorgehoben. Die Quantengravitation und die Quantenphysik Schwarzer Löcher zählten zu den zentralen Problemen der Grundlagenphysik. Die klassische Astrophysik und die numerische Relativitätstheorie untersuchten Systeme wie umeinander kreisende und verschmelzende Schwarze Löcher, bei denen die Lösung der Einsteinschen Gleichungen und die Vorhersage der Gravitationsstrahlung weit mehr erfordern, als Gleichungen blind in einen Computer einzugeben. Dieselben mathematischen Strukturen fanden auch in der Festkörperphysik und der Strömungsmechanik Anwendung. Ebenso bemerkenswert war der historische Wandel: Die allgemeine Relativitätstheorie entwickelte sich von einem vergleichsweise kleinen Spezialgebiet in den 1960er Jahren zu einer der zentralen Sprachen der theoretischen Physik, eng verbunden mit der Quantenfeldtheorie.
\lecturetimestamp{00:05:15}}
\end{classroomqa}

Dies ist die letzte Vorlesung, die sich ausschließlich mit Riemannscher Geometrie befasst. Von der nächsten Vorlesung an wird die Geometrie als Gravitation interpretiert. Die Verbindung ist bereits erkennbar. Ein echtes Gravitationsfeld von einem scheinbaren Feld zu unterscheiden, das durch beschleunigte oder krummlinige Koordinaten entsteht, ist mathematisch dasselbe Problem wie die Unterscheidung einer gekrümmten Metrik von einer flachen Metrik in ungünstigen Koordinaten.

\section{Die invariante Frage}

Eine zweidimensionale Geometrie ist flach, wenn die euklidische Geometrie intrinsisch gilt: Geraden, parallele Geraden, Dreiecksbeziehungen und Abstandsmessungen haben dann ihre euklidische Form. Es liegt nahe, Flachheit durch \(g_{mn}=\delta_{mn}\) zu definieren, doch diese Aussage hängt von den Koordinaten ab. Selbst eine gewöhnliche Ebene hat in Polarkoordinaten eine ortsabhängige Metrik. Die korrekte Aussage lautet, dass es für einen flachen Bereich Koordinaten gibt, in denen die Metrik im gesamten Bereich konstant und euklidisch ist.

Nehmen wir stattdessen an, die Metrik sei in beliebigen Koordinaten gegeben. Ein ungeeigneter Test wäre, Koordinatensysteme eines nach dem anderen durchzuprobieren, die Metrik jedes Mal zu transformieren und zu hoffen, schließlich eine kartesische Form zu finden. Da es unendlich viele Möglichkeiten gibt, ist das kein praktikables Kriterium.

Die richtige Strategie besteht darin, aus \(g_{mn}\) und seinen Ableitungen eine diagnostische Größe zu bilden. Gesucht ist ein Objekt \(R\) mit zwei entscheidenden Eigenschaften:
\[
\begin{aligned}
R=0
&\quad\Longleftrightarrow\quad
\text{locally flat geometry},\\
R\neq 0 \text{ at } P
&\quad\Longrightarrow\quad
\text{curvature at } P.
\end{aligned}
\]
Vor allem muss \(R\) ein Tensor sein. Ein Tensor, der in einem Koordinatensystem verschwindet, verschwindet in jedem Koordinatensystem; eine Suche nach passenden Koordinaten ist daher nicht nötig.

\begin{figure}[H]
\centering
\includegraphics[width=\linewidth]{lecture_03_figure_01.png}
\caption{Der gewünschte invariante Test: Eine Metrik muss nicht auf den ersten Blick \(\delta_{mn}\) sein; ein verschwindendes Krümmungsobjekt sollte dagegen eine lokal flache Geometrie kennzeichnen.}
\lectureframe{00:18:20}
\end{figure}

Die Einschränkung \emph{lokal} ist wichtig. Ein verschwindender Riemann-Tensor garantiert lokale euklidische Geometrie; globale Identifikationen können dennoch zu einer nichttrivialen Topologie führen. Ein flacher Zylinder oder ein flacher Torus ist beispielsweise lokal euklidisch, ohne global eine Ebene zu sein.

\section{Metrische Geometrie im Kleinen}

Es ist hilfreich, sich anzusehen, was sich mit der Dimension ändert. Ein nulldimensionaler Raum besteht nur aus einem Punkt. Ein zusammenhängender eindimensionaler Riemannscher Raum lässt sich lokal zu einer Geraden strecken; global kann er ein Intervall, eine Gerade oder ein geschlossener Weg sein. Ein Käfer, der auf eine geschlossene eindimensionale Bahn beschränkt ist, kann deren Umfang messen, doch eine Einbettung der Bahn in einen umgebenden Raum ändert ihre intrinsische Metrik nicht.

In zwei Dimensionen gibt es eine viel reichhaltigere intrinsische Geometrie: Eine Ebene, eine Kugel, eine unebene Erdoberfläche und ein Torus müssen weder lokal übereinstimmen noch dieselbe globale Topologie haben. Die Metrik erfasst ihre infinitesimalen Abstände:
\begin{equation}
 ds^2=g_{mn}(x)\,dx^m dx^n.
\label{eq:l3-line-element}
\end{equation}
Die Größe \(ds^2\) ist ein Skalar. Daher werden Änderungen von \(dx^m\) bei einer Koordinatentransformation genau durch die Transformation von \(g_{mn}\) ausgeglichen.

\begin{figure}[H]
\centering
\includegraphics[width=\linewidth]{lecture_03_figure_02.png}
\caption{Das invariante Linienelement mit einer ortsabhängigen Metrik.}
\lectureframe{00:19:00}
\end{figure}

Die Metrik ist nicht ausgeartet und besitzt daher eine Inverse:
\begin{equation}
 g^{mn}g_{nr}=\delta^m{}_r.
\label{eq:l3-inverse-metric}
\end{equation}
Der gemischte Kronecker-Tensor \(\delta^m{}_r\) ist die Identitätsabbildung. In der gewöhnlichen Riemannschen Geometrie ist die Metrik zudem positiv definit:
\begin{equation}
 g_{mn}u^m u^n>0
\qquad\text{for every nonzero }u^m.
\label{eq:l3-positive}
\end{equation}
Die Lorentzsche Raumzeit wird später auf diese Positivität verzichten, die Nichtausgeartetheit aber beibehalten.

\begin{figure}[H]
\centering
\includegraphics[width=\linewidth]{lecture_03_figure_03.png}
\caption{Die inverse Metrik und der Identitätstensor an der Tafel.}
\lectureframe{00:20:40}
\end{figure}

Die Metrik identifiziert Vektoren mit Kovektoren:
\begin{equation}
 V_m=g_{mn}V^n,
 \qquad
 V^m=g^{mn}V_n.
\label{eq:l3-raise-lower}
\end{equation}
Einen Index zu verschieben ist also eine Operation mit der Metrik und keine bloße Änderung der Schreibweise.

\begin{figure}[H]
\centering
\includegraphics[width=\linewidth]{lecture_03_figure_04.png}
\caption{Heben und Senken eines Vektorindex mit der Metrik und ihrer Inversen.}
\lectureframe{00:24:40}
\end{figure}

Sind alle Komponenten von \(g_{mn}\) in einem Bereich konstant, so kann eine konstante lineare Koordinatentransformation die Metrik dort diagonalisieren und normieren. Die Geometrie ist dann flach. Die interessante Information muss daher damit zusammenhängen, wie sich die Metrik mit dem Ort ändert.

\section{Normalkoordinaten: bis zur ersten Ordnung flach}

Wähle einen Punkt \(P\). Der Satz über Normalkoordinaten besagt, dass sich Koordinaten mit Ursprung in \(P\) konstruieren lassen, für die
\begin{align}
 g_{mn}(P)&=\delta_{mn}, \label{eq:l3-normal-metric}\\
 \partial_r g_{mn}(P)&=0. \label{eq:l3-normal-first}
\end{align}
In der Vorlesung werden sie Gaußsche Normalkoordinaten genannt. Geometrisch zeichnet man durch \(P\) Achsen und verlängert ihre Koordinatenlinien so geradlinig, wie es die Geometrie erlaubt. Was genau Geradlinigkeit bedeutet, wird sich als geodätische Bewegung erweisen.

\begin{figure}[H]
\centering
\includegraphics[width=\linewidth]{lecture_03_figure_05.png}
\caption{Normalkoordinaten, in denen die Metrik am ausgewählten Punkt euklidisch ist.}
\lectureframe{00:28:50}
\end{figure}

Ein solches Bezugssystem ist nicht eindeutig. Eine orthogonale Drehung bei \(P\) lässt \(g_{mn}(P)=\delta_{mn}\) unverändert. Entscheidend ist nicht die Orientierung der Achsen, sondern dass es ein Bezugssystem gibt, in dem sowohl \eqref{eq:l3-normal-metric} als auch \eqref{eq:l3-normal-first} gelten.

\begin{classroomqa}
\audiencequestion{Gelten diese Bedingungen nur am ausgewählten Punkt oder in einer ganzen Umgebung?}
\lecturerresponse{Sie gelten am ausgewählten Punkt \(P\). Man kann ein Normalkoordinatensystem um jeden Punkt zentrieren, aber im Allgemeinen lässt sich in einem gekrümmten Bereich kein einziges Koordinatensystem verwenden, in dem die Metrik überall euklidisch ist und alle ihre ersten Ableitungen verschwinden. Setzt man dieselben Koordinaten von \(P\) aus fort, treten höhere Ableitungen der Metrik zutage.
\lecturetimestamp{00:29:53}}
\end{classroomqa}

Bei \(P\) sind die zweiten Ableitungen im Allgemeinen unvermeidlich:
\begin{equation}
 \partial_r\partial_s g_{mn}(P)\neq 0
 \qquad\text{in general}.
\label{eq:l3-normal-second}
\end{equation}
Die Metrik enthält an einem Punkt also weder in nullter noch in erster Ordnung invariante Krümmungsinformation; die Koordinatenfreiheit kann diese Beiträge beseitigen. Die erste mögliche invariante Information tritt in zweiter Ordnung auf.

\begin{figure}[H]
\centering
\includegraphics[width=\linewidth]{lecture_03_figure_06.png}
\caption{Die Bedingungen für Normalkoordinaten: Die Metrik ist euklidisch, ihre ersten Ableitungen verschwinden, ihre zweiten Ableitungen müssen es nicht.}
\lectureframe{00:30:50}
\end{figure}

\subsection{Warum die Koordinatenfreiheit ausreicht}

Setze \(P\) als gemeinsamen Ursprung der Koordinaten \(x^m\) und \(y^m\). Nachdem eine lineare Transformation die Metrik am Ursprung in die gewünschte Form gebracht hat, betrachten wir den quadratischen Anteil einer weiteren Koordinatentransformation:
\begin{equation}
 x^m
 =y^m+C^m{}_{nr}y^n y^r+O(y^3),
 \qquad C^m{}_{nr}=C^m{}_{rn}.
\label{eq:l3-quadratic-transform}
\end{equation}
In vier Dimensionen hat das symmetrische Indexpaar \(nr\)
\[
\frac{4(4+1)}{2}=10
\]
unabhängige Möglichkeiten. Für jeden der vier Werte von \(m\) gibt es zehn Koeffizienten \(C^m{}_{nr}\), also insgesamt vierzig quadratische Koeffizienten.

Auch die symmetrische Metrik hat in vier Dimensionen zehn unabhängige Komponenten. Jede Komponente hat vier erste Ableitungen, sodass die Bedingungen \(\partial_r g_{mn}(P)=0\) vierzig Gleichungen umfassen. Diese übereinstimmende Anzahl zeigt, dass die quadratische Koordinatenfreiheit genau die richtige Größe hat, um die Bedingungen an die ersten Ableitungen durchzusetzen. Der Satz über Normalkoordinaten garantiert die Existenz; die Zählung ist eine nützliche strukturelle Kontrolle, aber kein vollständiger Beweis der Lösbarkeit eines beliebigen nichtlinearen Gleichungssystems.

\begin{figure}[H]
\centering
\includegraphics[width=\linewidth]{lecture_03_figure_07.png}
\caption{Die quadratische Koordinatentransformation und die zehn unabhängigen quadratischen Monome in vier Dimensionen.}
\lectureframe{00:33:50}
\end{figure}

\begin{classroomqa}
\audiencequestion{Ist \(x^m=y^m+C^m{}_{nr}y^n y^r+\cdots\) der Versuch einer Koordinatentransformation?}
\lecturerresponse{Ja. Die Koeffizienten \(C^m{}_{nr}\) werden so gewählt, dass die ersten Ableitungen der Metrik in den neuen Koordinaten am gemeinsamen Ursprung verschwinden. Der lineare Anteil hält die Jacobi-Matrix in der Nähe dieses Punktes invertierbar, während der quadratische Anteil die ersten Metrikableitungen anpasst. \lecturetimestamp{00:35:41}}
\end{classroomqa}

\begin{classroomqa}
\audiencequestion{Bedeutet es, dass eine glatte Geometrie lokal flach ist?}
\lecturerresponse{Es ist die entsprechende Aussage erster Ordnung. Auf hinreichend kleinen Skalen wird eine glatte Fläche durch ihre Tangentialebene angenähert, und Normalkoordinaten bringen dies durch \(g_{mn}(P)=\delta_{mn}\) und \(\partial_r g_{mn}(P)=0\) zum Ausdruck. Exakte Flachheit in einer Umgebung ist eine stärkere Forderung: Dort muss auch der Krümmungstensor verschwinden.
\lecturetimestamp{00:36:19}}
\end{classroomqa}

\begin{classroomqa}
\audiencequestion{Stellt die Gleichheit von vierzig Koeffizienten und vierzig Gleichungen wirklich sicher, dass sich die Gleichungen lösen lassen?}
\lecturerresponse{Die Zählung zeigt, welche Koordinatenfreiheit benötigt wird; der Satz garantiert die lokale Konstruktion für eine glatte, nicht ausgeartete Metrik. In der Vorlesung wird vorgeschlagen, den zweidimensionalen Fall als Übung direkt zu überprüfen. Das Argument ist nicht auf vier Dimensionen beschränkt: In \(d\) Dimensionen greifen dieselbe Symmetrie und Koordinatenfreiheit ineinander. \lecturetimestamp{00:36:45}}
\end{classroomqa}

\begin{classroomqa}
\audiencequestion{Warum gibt es in der Entwicklung keinen konstanten Term, und welcher Term ist von erster Ordnung?}
\lecturerresponse{Der Term \(y^m\) ist der Term erster Ordnung. Ein konstanter Term wurde weggelassen, weil die Ursprünge der beiden Koordinatensysteme bewusst so gewählt wurden, dass sie denselben Punkt darstellen; aus \(y=0\) folgt daher \(x=0\). Diese Wahl schränkt die Allgemeinheit nicht ein. \lecturetimestamp{00:38:52}}
\end{classroomqa}

\section{Warum die komponentenweise Ableitung scheitert}

Betrachten wir nun ein kovariantes Vektorfeld \(V_m(x)\). In geradlinigen kartesischen Koordinaten eines flachen Raums hat ein geometrisch konstanter Vektor konstante Komponenten. Man könnte daher versuchen, seine Ableitung als \(\partial_rV_m\) zu definieren. Das Problem zeigt sich, sobald sich das Koordinatengitter krümmt.

\begin{figure}[H]
\centering
\includegraphics[width=\linewidth]{lecture_03_figure_08.png}
\caption{Ein krummliniges Koordinatengitter, überlagert von einem lokal geradlinigen Bezugssystem.}
\lectureframe{00:40:10}
\end{figure}

Platzieren wir denselben geometrischen Vektor an benachbarten Punkten. Bezüglich eines festen kartesischen Bezugssystems ändern sich seine Komponenten nicht. Bezüglich einer krummlinigen Basis drehen sich dagegen die Basisvektoren mit der Bewegung, sodass sich die Komponentenliste ändert, obwohl der Vektor selbst unverändert bleibt.

\begin{figure}[H]
\centering
\includegraphics[width=\linewidth]{lecture_03_figure_09.png}
\caption{Derselbe Vektor hat an benachbarten Punkten unterschiedliche Komponenten, weil sich die lokale Koordinatenbasis gedreht hat.}
\lectureframe{00:42:30}
\end{figure}

Die naive Bedingung
\begin{equation}
 \partial_rV_m=0
\label{eq:l3-naive-derivative}
\end{equation}
kann also in einem Koordinatensystem gelten und in einem anderen nicht. Wäre \(\partial_rV_m\) ein Tensor vom Rang zwei, wäre das unmöglich: Sein Verschwinden in einem Bezugssystem würde sein Verschwinden in jedem Bezugssystem bedeuten. Gewöhnliche partielle Ableitungen von Vektorkomponenten sind folglich keine Komponenten eines Tensors.

\begin{figure}[H]
\centering
\includegraphics[width=\linewidth]{lecture_03_figure_10.png}
\caption{Die bewusst naive Bedingung für eine verschwindende Ableitung neben der sich ändernden krummlinigen Basis.}
\lectureframe{00:44:00}
\end{figure}

Wir benötigen eine Ableitung, die die scheinbare Änderung durch die Basis beseitigt und zugleich die echte Änderung des Vektorfelds beibehält.

\section{Definition der kovarianten Ableitung}

Bei der Ableitung in \(P\) werden nur Werte in einer beliebig kleinen Umgebung von \(P\) verglichen. Wir können daher zunächst Normalkoordinaten verwenden, die dort zentriert sind. Verschiebe einen Vektor an einem benachbarten Punkt so, dass sein Fußpunkt nach \(P\) zurückverlegt wird, und halte seine Normalkoordinatenkomponenten während dieser Verschiebung fest. Der Vergleich dieses transportierten Vektors mit dem bereits bei \(P\) befindlichen Vektor ergibt, geteilt durch den Abstand, die gewöhnliche Ableitung in diesem lokal geradlinigen Bezugssystem.

\begin{figure}[H]
\centering
\includegraphics[width=\linewidth]{lecture_03_figure_11.png}
\caption{Der lokale Vergleich: Vor der Differenzbildung wird der Vektor am benachbarten Punkt in Normalkoordinaten verschoben.}
\lectureframe{00:46:50}
\end{figure}

Nun drücken wir diese geometrisch definierte Ableitung in beliebigen Koordinaten aus. Es gibt zwei Beiträge: Die Komponenten \(V_m\) können sich ändern, und die Koordinatenbasis kann sich ändern. Die Indexstruktur legt die Form der Korrektur fest:
\begin{equation}
 \nabla_rV_m
 =\partial_rV_m-\Gamma^t{}_{rm}V_t.
\label{eq:l3-covector-derivative}
\end{equation}
Die Größen \(\Gamma^t{}_{rm}\) sind die Koeffizienten der Verbindung, auch Christoffelsymbole genannt. Das Minuszeichen ist für einen unteren Vektorindex passend. Im Zentrum \(P\) der Normalkoordinaten gilt
\begin{equation}
 \Gamma^t{}_{rm}(P)=0,
\qquad
 \nabla_rV_m(P)=\partial_rV_m(P).
\label{eq:l3-gamma-normal-zero}
\end{equation}

\begin{figure}[H]
\centering
\includegraphics[width=\linewidth]{lecture_03_figure_12.png}
\caption{Die kovariante Ableitung eines Kovektors: Änderung der Komponenten minus die durch die bewegte Basis hervorgerufene Änderung.}
\lectureframe{00:49:50}
\end{figure}

\begin{classroomqa}
\audiencequestion{Gilt \(\nabla_rV_m=\partial_rV_m-\Gamma^t{}_{rm}V_t\) in beliebigen Koordinaten, und wie vereinfacht sich der Ausdruck in Normalkoordinaten?}
\lecturerresponse{Die vollständige Formel gilt in beliebigen Koordinaten. Am ausgewählten Punkt in Normalkoordinaten verschwinden die ersten Ableitungen der Metrik und damit auch die Christoffelsymbole; es bleibt nur \(\partial_rV_m\). Außerhalb dieses Punktes sind die Verbindungkoeffizienten im Allgemeinen sogar in derselben Normalkoordinatenkarte ungleich null.
\lecturetimestamp{00:50:22}}
\end{classroomqa}

Die kovariante Ableitung heißt so, weil das gesamte Objekt \(\nabla_rV_m\) sich kovariant wie ein Tensor transformiert; der Name bezieht sich nicht bloß auf den unteren Ableitungsindex. Der Christoffelterm behebt genau das nichttensorielle Transformationsverhalten der partiellen Ableitung.

\begin{figure}[H]
\centering
\includegraphics[width=\linewidth]{lecture_03_figure_13.png}
\caption{Die vollständige Kovektor-Ableitung mit ihrer Verbindungskorrektur.}
\lectureframe{00:51:30}
\end{figure}

\begin{classroomqa}
\audiencequestion{Hat das etwas mit der konvektiven Ableitung in der Strömungsmechanik zu tun?}
\lecturerresponse{Die Vorlesung bejaht vorsichtig einen Zusammenhang: Auch die Strömungsmechanik unterscheidet zwischen einer Änderung, die von einer Strömung mitgeführt wird, und einer Änderung an einem festen Koordinatenpunkt; außerdem treten kovariante Ableitungen in Strömungsgleichungen auf gekrümmten Koordinatensystemen oder in gekrümmten Räumen auf. Die Operatoren sind nicht per Definition identisch, aber beide korrigieren eine naive Ableitung um die Struktur, die sich mit der Beschreibung bewegt.
\lecturetimestamp{00:52:31}}
\end{classroomqa}

Für die hier verwendete torsionsfreie Verbindung gilt
\begin{equation}
 \Gamma^t{}_{rm}=\Gamma^t{}_{mr}.
\label{eq:l3-torsion-free}
\end{equation}
Allgemeinere Geometrien können Torsion besitzen und müssen diese Symmetrie nicht erfüllen. Die hier betrachtete Verbindung ist diejenige, die für die übliche metrische Formulierung der allgemeinen Relativitätstheorie maßgeblich ist.

\begin{figure}[H]
\centering
\includegraphics[width=\linewidth]{lecture_03_figure_14.png}
\caption{Die Christoffelsymbole als symmetrische Verbindungkoeffizienten der torsionsfreien Geometrie.}
\lectureframe{00:54:30}
\end{figure}

Es gibt eine nützliche physikalische Analogie. Normalkoordinaten entsprechen dem Bezugssystem eines frei fallenden Beobachters an einem Ereignis: Dort lässt sich die Koordinatenkorrektur zum Verschwinden bringen. Beliebige krummlinige Koordinaten ähneln einem beschleunigten Aufzug, in dem scheinbare Kräfte auftreten, weil sich das Bezugssystem selbst ändert. Die Analogie gilt lokal; die Krümmung wird erklären, warum ein einziges frei fallendes Bezugssystem die Gravitation nicht in einem ausgedehnten Bereich beseitigen kann.

\subsection{Eine Korrektur für jeden Tensorindex}

Bei einem kovarianten Tensor zweiten Ranges \(T_{mn}\) erhält jeder untere Index seine eigene Verbindungskorrektur:
\begin{equation}
 \nabla_sT_{mn}
 =\partial_sT_{mn}
 -\Gamma^t{}_{sm}T_{tn}
 -\Gamma^t{}_{sn}T_{mt}.
\label{eq:l3-rank-two-derivative}
\end{equation}
Zur Erinnerung an die Regel hält man beim Ableiten des einen Index den anderen als Zuschauer fest und wiederholt den Vorgang dann für den zweiten Index. Kontravariante Indizes erhalten entsprechende Terme mit umgekehrtem Vorzeichen.

\begin{figure}[H]
\centering
\includegraphics[width=\linewidth]{lecture_03_figure_15.png}
\caption{Die kovariante Ableitung eines kovarianten Tensors zweiten Ranges mit einer Christoffelkorrektur für jeden unteren Index.}
\lectureframe{01:01:50}
\end{figure}

\begin{classroomqa}
\audiencequestion{Bezieht sich das Symbol \(\partial_r\) nur auf Normalkoordinaten, und gilt die Tensorformel nur in diesem Bezugssystem?}
\lecturerresponse{Nein. Die Schreibweise \(\partial_r=\partial/\partial x^r\) ist in jedem Koordinatensystem sinnvoll, und die vollständige Formel der kovarianten Ableitung gilt in beliebigen Koordinaten. Nur am ausgewählten Punkt eines Normalkoordinatensystems verschwinden die Christoffelkorrekturen. \lecturetimestamp{00:57:22}}
\end{classroomqa}

\begin{classroomqa}
\audiencequestion{Wenn die kovariante Ableitung in Normalkoordinaten auf die partielle Ableitung zurückgeführt wird, gilt das nur am ausgewählten Punkt? Sind die Christoffelsymbole Tensoren?}
\lecturerresponse{Ja, die Vereinfachung gilt punktweise. Die Christoffelsymbole sind keine Tensoren: In Normalkoordinaten können sie bei \(P\) verschwinden und am selben Punkt in einem anderen Koordinatensystem ungleich null sein. Gerade ihr nichttensorielles Transformationsverhalten ermöglicht es, dass sich die Kombination in \eqref{eq:l3-covector-derivative} wie ein Tensor transformiert.
\lecturetimestamp{00:58:13}}
\end{classroomqa}

\begin{classroomqa}
\audiencequestion{Wozu wird die kovariante Ableitung konkret verwendet?}
\lecturerresponse{Für Feldgleichungen. Gravitationsfelder ändern sich von Punkt zu Punkt, daher müssen ihre Gleichungen Ableitungen enthalten. Zugleich sollen die Gleichungen in jedem Koordinatensystem denselben Inhalt haben. Die kovariante Ableitung bildet aus Tensoren wieder Tensoren und ermöglicht so Differentialgleichungen, deren Gültigkeit nicht von einem besonderen Bezugssystem abhängt. \lecturetimestamp{01:02:35}}
\end{classroomqa}

\begin{classroomqa}
\audiencequestion{Kann man die kovariante Ableitung mit einem Koordinatendifferential multiplizieren, um die Änderung zwischen benachbarten Punkten zu beschreiben?}
\lecturerresponse{Ja. Entlang einer infinitesimalen Verschiebung \(dx^s\) ist die kovariante Änderung \(dx^s\nabla_sT\). Der partielle Term misst die Änderung der Komponenten, während die Verbindungsterme die Änderung der Koordinatenbasis berücksichtigen. Auch im flachen Raum in Polarkoordinaten muss man diese Unterscheidung machen. \lecturetimestamp{01:03:22}}
\end{classroomqa}

\section{Metrische Verträglichkeit bestimmt die Verbindung}

Wenden wir die Regel für Tensoren zweiten Ranges auf die Metrik an. Bei \(P\) gilt in Normalkoordinaten \(\partial_sg_{mn}=0\) und \(\Gamma^t{}_{sm}=0\), also \(\nabla_sg_{mn}=0\). Da die kovariante Ableitung ein Tensor ist, gilt dieses punktweise Ergebnis in jedem Koordinatensystem:
\begin{equation}
 \nabla_sg_{mn}=0.
\label{eq:l3-metric-compatibility}
\end{equation}
Dies ist die metrische Verträglichkeit. Sie besagt, dass der Paralleltransport die Metrik und damit Längen und Winkel erhält.

\begin{figure}[H]
\centering
\includegraphics[width=\linewidth]{lecture_03_figure_16.png}
\caption{Die Gleichung der metrischen Verträglichkeit leitet die Bestimmung der Verbindung aus der Metrik ein.}
\lectureframe{01:06:20}
\end{figure}

Durch Ausmultiplizieren von \eqref{eq:l3-metric-compatibility} ergibt sich
\begin{equation}
 0
 =\partial_sg_{mn}
 -\Gamma^t{}_{sm}g_{tn}
 -\Gamma^t{}_{sn}g_{mt}.
\label{eq:l3-compat-expanded}
\end{equation}
Schreiben wir noch zwei Gleichungen mit zyklisch vertauschten Ableitungs- und Metrikindizes:
\begin{align}
 0
 &=\partial_mg_{sn}
 -\Gamma^t{}_{ms}g_{tn}
 -\Gamma^t{}_{mn}g_{st},
\label{eq:l3-compat-m}\\
 0
 &=\partial_ng_{sm}
 -\Gamma^t{}_{ns}g_{tm}
 -\Gamma^t{}_{nm}g_{st}.
\label{eq:l3-compat-n}
\end{align}

\begin{figure}[H]
\centering
\includegraphics[width=\linewidth]{lecture_03_figure_17.png}
\caption{Drei Gleichungen der metrischen Verträglichkeit mit vertauschten Indizes.}
\lectureframe{01:10:20}
\end{figure}

Addiere \eqref{eq:l3-compat-m} und \eqref{eq:l3-compat-n} und subtrahiere anschließend \eqref{eq:l3-compat-expanded}. Aufgrund der Torsionsfreiheit und der Symmetrie \eqref{eq:l3-torsion-free} heben sich die unerwünschten Verbindungsterme auf:
\begin{equation}
 \partial_mg_{sn}+\partial_ng_{sm}-\partial_sg_{mn}
 =2\Gamma^t{}_{mn}g_{st}.
\label{eq:l3-isolate-gamma-lowered}
\end{equation}

\begin{figure}[H]
\centering
\includegraphics[width=\linewidth]{lecture_03_figure_18.png}
\caption{Durch Addition zweier Verträglichkeitsgleichungen und Subtraktion der dritten heben sich die unerwünschten Christoffelterme auf.}
\lectureframe{01:12:10}
\end{figure}

Die Multiplikation mit \(g^{rs}\) isoliert die Verbindung:
\begin{equation}
 \boxed{
 \Gamma^r{}_{mn}
 =\frac12 g^{rs}
 \left(
 \partial_mg_{sn}
 +\partial_ng_{sm}
 -\partial_sg_{mn}
 \right)}
\label{eq:l3-christoffel}
\end{equation}
Die Levi-Civita-Verbindung ist somit aus der Metrik und ihren ersten Ableitungen aufgebaut. Damit ist erklärt, warum die Christoffelsymbole im Zentrum der Normalkoordinaten verschwinden.

\begin{figure}[H]
\centering
\includegraphics[width=\linewidth]{lecture_03_figure_19.png}
\caption{Die inverse Metrik isoliert das Christoffelsymbol aus der abgesenkten Verbindungsgleichung.}
\lectureframe{01:16:10}
\end{figure}

\begin{classroomqa}
\audiencequestion{Warum wurden alle drei Gleichungen der metrischen Verträglichkeit gleich null gesetzt?}
\lecturerresponse{Am ausgewählten Punkt in Normalkoordinaten verschwinden die gewöhnlichen Ableitungen der Metrik ebenso wie die Christoffelsymbole. Daher ist dort auch die kovariante Ableitung der Metrik null. Um die Aussage an einem anderen Punkt zu zeigen, zentriert man dort ein Normalkoordinatensystem und wiederholt die Rechnung. Da \(\nabla_sg_{mn}\) ein Tensor ist, ist sein Verschwinden koordinatenunabhängig. \lecturetimestamp{01:14:27}}
\end{classroomqa}

\begin{classroomqa}
\audiencequestion{Welche Dimensionen haben die Metrik und das Christoffelsymbolarray, und warum gibt es in vier Dimensionen vierzig unabhängige Christoffelkomponenten?}
\lecturerresponse{In \(d\) Dimensionen ist \(g_{mn}\) eine symmetrische Matrix der Größe \(d\times d\). Die Verbindung hat drei Indizes. Für jeden festen oberen Index \(r\) gibt es für das symmetrische untere Indexpaar \(mn\) genau \(d(d+1)/2\) Möglichkeiten. Eine torsionsfreie Verbindung hat somit \(d^2(d+1)/2\) unabhängige Komponenten; für \(d=4\) sind es vierzig.
\lecturetimestamp{01:16:18}}
\end{classroomqa}

\begin{figure}[H]
\centering
\includegraphics[width=\linewidth]{lecture_03_figure_20.png}
\caption{Die Tafel während der Diskussion über die Komponentenanzahl: zehn Möglichkeiten für das symmetrische untere Indexpaar und vier Möglichkeiten für den oberen Index.}
\lectureframe{01:16:30}
\end{figure}

\begin{classroomqa}
\audiencequestion{Welche Metrik tritt in \eqref{eq:l3-christoffel} auf? Muss man an jedem Punkt eine andere Metrik in Normalkoordinaten verwenden?}
\lecturerresponse{Die Formel verwendet die Komponenten der Metrik in den globalen oder lokalen Koordinaten, die man für die Rechnung gewählt hat. Normalkoordinaten sind lediglich das lokale Hilfsmittel, mit dem die Regel hergeleitet wird: Man wechselt an einem Punkt in ein Normalkoordinatensystem, bildet dort die Ableitung, transformiert das Ergebnis zurück und verwirft anschließend das Hilfsbezugssystem. Sobald \eqref{eq:l3-christoffel} bekannt ist, wertet man es direkt in den interessierenden beliebigen Koordinaten aus. \lecturetimestamp{01:17:21}}
\end{classroomqa}

\begin{classroomqa}
\audiencequestion{Sagt entweder \(\nabla_sg_{mn}=0\) oder die Christoffelformel aus, dass der Raum flach ist?}
\lecturerresponse{Nein. Die metrische Verträglichkeit gilt für jede Levi-Civita-Verbindung, ob der Raum gekrümmt oder flach ist. Christoffelsymbole können auch auf einer flachen Ebene ungleich null sein, wenn man Polarkoordinaten verwendet. Verschwinden alle \(\Gamma^r{}_{mn}\) in einem Bereich in einer Karte, so ist diese Karte dort kartesisch und der Bereich flach; ihr bloßes Nichtverschwinden weist jedoch keine Krümmung nach. Ein aus Ableitungen und Produkten von \(\Gamma\) gebildeter Tensor ist weiterhin erforderlich. \lecturetimestamp{01:19:27}}
\end{classroomqa}

\section{Ein Kegel zeigt, was die Verbindung nicht leisten kann}

Betrachten wir einen abgerundeten Kegel. Abseits seiner Spitze kann er durch einen gewöhnlichen Kegel aus Papier angenähert werden. Man schneidet ein Kreissegment aus einer flachen Scheibe heraus und klebt die freigelegten Ränder zusammen. Das verbleibende Papier wurde an keiner Stelle gedehnt; lokal ist jeder Punkt abseits der Naht und der Spitze weiterhin flach. Der fehlende Sektor zeigt sich im Defizitwinkel \(\Delta\).

\begin{figure}[H]
\centering
\includegraphics[width=\linewidth]{lecture_03_figure_21.png}
\caption{Der Kegel als Beispiel einer Geometrie, die abseits der Spitze lokal flach ist, deren Krümmung aber an der Spitze konzentriert ist.}
\lectureframe{01:23:00}
\end{figure}

\begin{figure}[H]
\centering
\includegraphics[width=\linewidth]{lecture_03_figure_22.png}
\caption{Durch Herausschneiden eines Kreissektors aus einer Scheibe und Identifizieren ihrer Ränder entsteht der Defizitwinkel des Kegels.}
\lectureframe{01:24:20}
\end{figure}

Die Konstruktion aus Papier macht den Unterschied anschaulich. Auf dem aufgeschnittenen Blatt transportiert man einen Vektor parallel in der gewöhnlichen Ebene. Nachdem die Schnittkanten zusammengefügt wurden, kehrt ein Weg, der die Spitze umrundet, zum Ausgangspunkt zurück; der transportierte Vektor ist dabei gegenüber seiner ursprünglichen Richtung gedreht.

\begin{figure}[H]
\centering
\includegraphics[width=0.78\linewidth]{lecture_03_figure_23.png}
\caption{Die Papierkegel-Demonstration: Aus einem flachen Blatt, aus dem ein Keil herausgeschnitten wurde, entsteht ein Kegel, wenn man die Schnittkanten zusammenfügt.}
\lectureframe{01:25:10}
\end{figure}

\begin{figure}[H]
\centering
\includegraphics[width=\linewidth]{lecture_03_figure_24.png}
\caption{Vektoren auf dem ausgeschnittenen Keil, bevor die Kanten zusammengefügt werden.}
\lectureframe{01:27:40}
\end{figure}

Diese Drehung ist Holonomie. Sie wird nicht durch eine ungünstige Wahl der Koordinaten entlang eines einzelnen Abschnitts verursacht. Jeder kleine Abschnitt lässt sich in einem lokal flachen Bezugssystem behandeln. Der Effekt tritt erst auf, wenn man diese lokalen Vergleiche entlang einer geschlossenen Schleife zusammensetzt. Beim idealen Kegel entspricht die resultierende Drehung, bis auf die Orientierung, dem Defizitwinkel.

\begin{figure}[H]
\centering
\includegraphics[width=\linewidth]{lecture_03_figure_25.png}
\caption{Beim Paralleltransport um die Kegelspitze kehrt ein Vektor mit gedrehter Richtung zurück: Die Holonomie erfasst die eingeschlossene Krümmung.}
\lectureframe{01:29:10}
\end{figure}

\section{Zwei Wege um eine infinitesimale Schleife}

Ersetze die endliche Schleife durch ein winziges Koordinatenrechteck. Beginne an einer Ecke und transportiere zuerst in \(r\)-Richtung und dann in \(s\)-Richtung. Vergleiche dieses Ergebnis mit dem Transport zuerst in \(s\)-Richtung und dann in \(r\)-Richtung. In flachen kartesischen Koordinaten stimmen die beiden Wege überein. In einer gekrümmten Geometrie ist das im Allgemeinen nicht der Fall.

\begin{figure}[H]
\centering
\includegraphics[width=\linewidth]{lecture_03_figure_26.png}
\caption{Zwei Reihenfolgen des infinitesimalen Transports um ein Koordinatenrechteck.}
\lectureframe{01:31:40}
\end{figure}

Das algebraische Maß dafür ist der Kommutator kovarianter Ableitungen:
\begin{equation}
 \left(\nabla_s\nabla_r-\nabla_r\nabla_s\right)V_n.
\label{eq:l3-commutator}
\end{equation}
Gewöhnliche gemischte partielle Ableitungen vertauschen auf glatten Funktionen. Dass dies hier nicht gilt, liegt daran, dass neben den Vektorkomponenten auch die Terme mit der Verbindung abgeleitet werden.

\begin{figure}[H]
\centering
\includegraphics[width=\linewidth]{lecture_03_figure_27.png}
\caption{Die infinitesimale Schleife, dargestellt durch die Differenz zwischen den beiden Reihenfolgen der kovarianten Ableitung.}
\lectureframe{01:33:20}
\end{figure}

Zur Berechnung betrachte man zunächst
\[
 T_{rn}=\nabla_rV_n
\]
als kovarianten Tensor zweiter Stufe. Verwende dann
\eqref{eq:l3-rank-two-derivative}, um \(T_{rn}\) in \(s\)-Richtung abzuleiten. Wiederhole das mit vertauschten \(r\) und \(s\) und bilde die Differenz. Die Rechnung ist lang, aber rein mechanisch: Man entwickelt die Terme, wendet die Produktregel an und kürzt die Terme, die zweite partielle Ableitungen von \(V_n\) enthalten.

\begin{figure}[H]
\centering
\includegraphics[width=\linewidth]{lecture_03_figure_28.png}
\caption{Beginn der doppelten kovarianten Ableitung, wobei \(\nabla_rV_n\) als Tensor zweiter Stufe behandelt wird.}
\lectureframe{01:35:40}
\end{figure}

Mit der an der Tafel verwendeten Konvention ergibt sich
\begin{equation}
 \left(\nabla_s\nabla_r-\nabla_r\nabla_s\right)V_n
 =R_{srn}{}^tV_t.
\label{eq:l3-riemann-commutator}
\end{equation}

\begin{figure}[H]
\centering
\includegraphics[width=\linewidth]{lecture_03_figure_29.png}
\caption{Der Kommutator kovarianter Ableitungen definiert die Wirkung des Krümmungstensors auf einen Kovektor.}
\lectureframe{01:38:30}
\end{figure}

Für diese Reihenfolge im Kommutator gilt:
\begin{equation}
\begin{split}
 R_{srn}{}^t
={}&
 \partial_r\Gamma^t{}_{sn}
 -\partial_s\Gamma^t{}_{rn}\\
&+\Gamma^p{}_{sn}\Gamma^t{}_{pr}
 -\Gamma^p{}_{rn}\Gamma^t{}_{ps}.
\end{split}
\label{eq:l3-riemann}
\end{equation}
Andere Texte verwenden möglicherweise das umgekehrte Vorzeichen oder ordnen die Indizes anders an; die Konvention ist durch die Angabe sowohl von \eqref{eq:l3-riemann-commutator} als auch von \eqref{eq:l3-riemann} festgelegt.

\begin{figure}[H]
\centering
\includegraphics[width=\linewidth]{lecture_03_figure_30.png}
\caption{Der Riemann-Tensor: zwei Ableitungen der Verbindung und zwei quadratische Verbindungsterme.}
\lectureframe{01:39:40}
\end{figure}

Gleichung \eqref{eq:l3-christoffel} zeigt, weshalb die Krümmung erstmals in zweiter Ordnung der Metrik auftritt. Jedes \(\Gamma\) enthält eine erste Ableitung von \(g\). Die Ableitungsterme in \eqref{eq:l3-riemann} enthalten daher zweite Ableitungen von \(g\), während die quadratischen Terme Produkte erster Ableitungen enthalten. In Normalkoordinaten bei \(P\) verschwindet die Verbindung selbst, ihre Ableitungen im Allgemeinen jedoch nicht, sodass die Krümmung erhalten bleibt:
\[
\begin{aligned}
 \Gamma^t{}_{mn}(P)&=0,\\
 R_{srn}{}^t(P)
 &=\partial_r\Gamma^t{}_{sn}(P)
 -\partial_s\Gamma^t{}_{rn}(P).
\end{aligned}
\]
Es ist wichtig, die Ableitungsordnung im Blick zu behalten: Der Riemann-Tensor verwendet erste Ableitungen von \(\Gamma\), die in der Metrik Ableitungen zweiter Ordnung sind.

Dies ist das gesuchte Diagnosekriterium. Verschwinden alle Komponenten des Riemann-Tensors in einem Gebiet, so ist dieses Gebiet lokal flach. Ist an einem Punkt irgendeine Komponente ungleich null, lässt sich die Krümmung dort durch keinen Koordinatenwechsel beseitigen. Ein Computeralgebrasystem kann diese Komponenten ausgehend von einer analytisch gegebenen Metrik prüfen; der Begriff selbst bezeichnet jedoch schlicht das Ausbleiben des Schließens beim Paralleltransport um eine infinitesimale Schleife.

\section{Krümmung als Gezeitenwirkung der Gravitation}

Das geometrische Diagnosekriterium hat eine unmittelbare physikalische Bedeutung. Man stelle sich einen biegsamen Gegenstand aus miteinander verbundenen Dreiecken vor. In einem flachen Gebiet lässt er sich bewegen, ohne dass Spannungen auftreten. Bringt man ihn in ein Gebiet, in dem die intrinsische Metrik gekrümmt ist, können die verschiedenen Teile des Verbunds ihre bevorzugten Abstände und Winkel nicht alle beibehalten. Der Gegenstand muss sich dehnen, stauchen oder scheren.

\begin{figure}[H]
\centering
\includegraphics[width=\linewidth]{lecture_03_figure_31.png}
\caption{Ein verbundener Gegenstand, der aus einem flachen Gebiet in eine gekrümmte Ausbuchtung bewegt wird, muss sich verformen, um der lokalen Metrik zu folgen.}
\lectureframe{01:44:10}
\end{figure}

Diese Verformung ist das geometrische Gegenstück zur Gezeitenkraft. In einem Gravitationsfeld können benachbarte frei fallende Teilchen relativ zueinander beschleunigt werden: In einer Richtung werden sie möglicherweise auseinandergezogen, in einer anderen zusammengedrückt. Der Riemann-Tensor bestimmt diese Relativbeschleunigung. In der später eingeführten Sprache der geodätischen Abweichung gilt für den Tangentialvektor \(u^\mu\) und den Abstandsvektor \(\xi^\mu\):
\begin{equation}
 \frac{D^2\xi^\mu}{D\tau^2}
 =-R^\mu{}_{\nu\rho\sigma}
 u^\nu\xi^\rho u^\sigma,
\label{eq:l3-geodesic-deviation-preview}
\end{equation}
wobei das Gesamtvorzeichen von der gewählten Krümmungskonvention abhängt. Die Gleichung wird hier nicht hergeleitet, präzisiert aber die Aussage der Vorlesung, dass Krümmung und Gezeitenspannung buchstäblich dieselbe lokale Physik sind.

Ein ideales, vollkommen homogenes Gravitationsfeld weist keine Gezeitenvariation auf und lässt sich lokal beseitigen, indem man in ein gleichmäßig beschleunigtes Bezugssystem wechselt. Reale Felder, die von isolierten Massen erzeugt werden, sind nicht vollkommen homogen. Das Gravitationsfeld der Erde ist annähernd radial gerichtet und ändert sich mit dem Ort, sodass benachbarte Fallbahnen konvergieren oder auseinanderlaufen.

\begin{figure}[H]
\centering
\includegraphics[width=\linewidth]{lecture_03_figure_32.png}
\caption{Ein radiales Gravitationsfeld variiert über einen ausgedehnten Körper hinweg und erzeugt so die Relativbeschleunigung, die als Gezeitenwirkung der Gravitation bezeichnet wird.}
\lectureframe{01:48:40}
\end{figure}

\begin{classroomqa}
\audiencequestion{Woran erkennt man ein echtes Gravitationsfeld im Unterschied zu dem scheinbaren Feld in einem beschleunigten Bezugssystem?}
\lecturerresponse{Ein gleichmäßig beschleunigtes Bezugssystem kann bei lokalen Experimenten ein homogenes Gravitationsfeld nachahmen, erzeugt aber keine Gezeiten-Relativbeschleunigung. Ein von realer Materie erzeugtes Feld ändert sich von Ort zu Ort und verursacht daher Gezeitenkräfte. Diese Kräfte sind ein invariantes Indiz dafür, dass der Effekt nicht bloß auf einem Koordinatenwechsel beruht. Ein vollkommen homogenes reales Feld ist eine Idealisierung; nahe genug an der Oberfläche eines großen Körpers ist das tatsächliche Feld nur näherungsweise homogen. \lecturetimestamp{01:46:54}}
\end{classroomqa}

\begin{classroomqa}
\audiencequestion{Warum spielt die Größe des fallenden Körpers für Gezeiteneffekte eine Rolle?}
\lecturerresponse{Krümmung beschreibt den Vergleich benachbarter Orte. Ein Körper, der viel kleiner ist als die Längenskala, über die sich das Feld ändert, erfährt an all seinen Teilen nahezu dieselbe Beschleunigung und verformt sich nur geringfügig. Ein ausgedehnter Körper erfährt größere Unterschiede. Ein frei fallendes Bakterium in Erdnähe spürt vernachlässigbare Gezeitenspannungen, während der in der Vorlesung ausgedachte Mensch mit einer Ausdehnung von zweitausend Meilen stark verformt würde. Die Krümmung hat einen lokalen Wert, doch die aufsummierte Verformung hängt vom verglichenen Abstand ab.
\lecturetimestamp{01:49:01}}
\end{classroomqa}

\begin{classroomqa}
\audiencequestion{Wenn man zur Feststellung der Flachheit die Krümmung an jedem Punkt prüfen muss, ist die Aufgabe dann immer noch unendlich?}
\lecturerresponse{Bei einer analytisch gegebenen Metrik sind ihre Ableitungen und die Krümmungskomponenten analytische Ausdrücke. Man berechnet sie einmal als Funktionen des Ortes und kann oft erkennen, ob sie identisch verschwinden. Liegt die Metrik nur als numerische Tabelle vor, muss man die Krümmung im gesamten abgetasteten Gebiet abschätzen und untersuchen, was tatsächlich aufwendig sein kann. Dennoch ist das wesentlich besser definiert, als an jedem Punkt jedes denkbare Koordinatensystem zu durchsuchen.
\lecturetimestamp{01:50:37}}
\end{classroomqa}

\section{Was bisher gezeigt wurde}

Die Argumentation lässt sich nun als zusammenhängende Kette lesen:
\begin{enumerate}
 \item An jedem gewählten Punkt lassen sich Koordinaten so wählen, dass \(g_{mn}=\delta_{mn}\) und \(\partial_rg_{mn}=0\) gelten.
 \item Gewöhnliche Ableitungen von Tensor-Komponenten transformieren sich nicht wie Tensoren, weil die Koordinatenbasen von Punkt zu Punkt variieren.
 \item Die Verbindung korrigiert diese Variation der Basis und definiert die kovariante Ableitung.
 \item Metrikverträglichkeit und verschwindende Torsion bestimmen die Verbindung eindeutig aus der Metrik.
 \item Zwei kovariante Ableitungen kommutieren nicht; ihre Differenz wird durch den Riemann-Tensor beschrieben. Dieses Nichtkommutieren misst die Holonomie, die Metrikvariation zweiter Ordnung und die Gezeitenwirkung der Gravitation.
\end{enumerate}

Damit wurde das angekündigte Diagnosekriterium für Krümmung hergeleitet, ohne ein ausgezeichnetes globales Koordinatensystem zu wählen. Als Nächstes wechseln wir von der positiv definiten riemannschen Geometrie zur lorentzschen Raumzeit, wenden diese Werkzeuge auf die Schwarzschild-Geometrie eines Sterns oder Schwarzen Lochs an und kehren zu den Geodäten zurück, die am Ende dieser Vorlesung zunächst zurückgestellt wurden.